\section{CLIP 与 SigLIP：先学会“图像和文字是否匹配”}

本节的问题是：没有密集人工类别标签，怎样利用互联网规模的 image-caption pairs？CLIP 用双 encoder 将图像和文本映射到共享 embedding space，并通过 batch 内 contrastive learning 拉近匹配 pair、推远不匹配 pair。

\subsection{ViT：把图像切成 patch tokens}

为了复用 Transformer，ViT 先把二维图像改写成序列。图像 $I\in\mathbb{R}^{H\times W\times C}$ 被切成 $P\times P$ patches。Patch 数量为

$$
N=\frac{H}{P}\cdot\frac{W}{P}.
$$

每个 patch 展平后经过线性 projection，加入位置编码，再送入 Transformer encoder。较小 $P$ 保留更多细节，却让 attention cost 随 $N^2$ 增长。

\begin{figure}[H]
\centering
\includegraphics[width=0.78\textwidth]{vit.png}
\caption{Vision Transformer：patchify、linear projection、position embedding 与 Transformer encoder。}
\figsource{Stanford CS336 2026 Lecture 17 官方可执行讲义，\texttt{images/vit.png}。}
\end{figure}

\begin{knowledgebox}{读图：ViT 把视觉问题变成序列问题}
Patch embedding 扮演 word embedding 的角色，position embedding 保留空间位置，CLS 或 pooled representation 汇总全图。它让视觉可以复用 Transformer，但也把 resolution 直接转化为 sequence length。
\end{knowledgebox}

\subsection{CLIP objective}

接下来需要让视觉序列与语言语义对齐。给定 batch 中 $B$ 个 image-text pairs，图像向量 $u_i$、文本向量 $v_j$ 归一化后计算相似度 $s_{ij}=u_i^\top v_j/\tau$。CLIP 同时优化 image-to-text 与 text-to-image cross entropy：

$$
\mathcal{L}_{\mathrm{CLIP}}
=\frac{1}{2B}\sum_{i=1}^{B}
\left[-\log\frac{e^{s_{ii}}}{\sum_j e^{s_{ij}}}
-\log\frac{e^{s_{ii}}}{\sum_j e^{s_{ji}}}\right].
$$

$\tau$ 是 temperature；batch 中其他 pairs 充当 negatives。

\begin{figure}[H]
\centering
\includegraphics[width=0.7\textwidth]{clip-code.png}
\caption{CLIP objective 的伪代码：归一化双 encoder embeddings，构造全 batch logits，并对两个方向做 cross entropy。}
\figsource{Stanford CS336 2026 Lecture 17 官方可执行讲义，\texttt{images/clip-code.png}。}
\end{figure}

\begin{knowledgebox}{读图：代码中的矩阵 shape 就是 contrastive objective}
图像与文本 encoder 分别输出 $[B,d]$，归一化后矩阵乘法得到 $[B,B]$ logits；对角线是 aligned pairs，其余位置都是 batch negatives。随后分别沿 image-to-text 与 text-to-image 方向做 cross entropy。这里最该注意的是 negatives 来自当前 global batch，所以代码看似短，真正的系统成本却在大 batch、跨设备 embedding 收集和 $B^2$ score matrix。
\end{knowledgebox}

\begin{knowledgebox}{课堂提示：CLIP 的 batch size 是 objective 的一部分}
老师指出，CLIP 典型 batch 可到数万；batch 太小不仅降低吞吐，还会让 negatives 数量不足，直接改变学习问题。更麻烦的是 softmax 分母跨越整个 batch，使设备之间无法像普通 LM sequences 那样完全独立计算。SigLIP 的动机正是解除这种 global normalization 耦合，而不是单纯换一个公式名字。
\end{knowledgebox}

\begin{figure}[H]
\centering
\includegraphics[width=0.84\textwidth]{clip.png}
\caption{CLIP：双 encoder 与 batch-wise contrastive matching。}
\figsource{Stanford CS336 2026 Lecture 17 官方可执行讲义，\texttt{images/clip.png}。}
\end{figure}

\begin{knowledgebox}{读图：大 batch 为什么重要}
每个 batch 提供 $B^2$ 个 pair scores，其中只有 $B$ 个正例。更大 batch 给出更多 negatives，通常提高 representation quality；但也需要跨设备 all-gather embeddings，并使 global softmax 成为通信与数值稳定性问题。
\end{knowledgebox}

\begin{figure}[H]
\centering
\includegraphics[width=0.82\textwidth]{clip-efficiency.png}
\caption{CLIP 训练效率与 batch/negative scaling 的关系。}
\figsource{Stanford CS336 2026 Lecture 17 官方可执行讲义，\texttt{images/clip-efficiency.png}。}
\end{figure}

\subsection{SigLIP：把 global softmax 换成 pairwise sigmoid}

进一步，SigLIP 取消 batch-wide softmax，把每个 image-text pair 变成 binary classification：匹配 pair 标签为 $+1$，不匹配为 $-1$。Loss 可写成

$$
\mathcal{L}_{\mathrm{SigLIP}}
=\frac{1}{B^2}\sum_{i,j}\log\left(1+\exp\left[-z_{ij}(u_i^\top v_j/\tau+b)\right]\right).
$$

$z_{ij}\in\{+1,-1\}$ 表示是否匹配，$b$ 是 bias。它避免 global softmax normalization，更容易在多设备上做 local pair computation。

\begin{figure}[H]
\centering
\includegraphics[width=0.84\textwidth]{siglip-code.png}
\caption{SigLIP sigmoid loss 伪实现：每个 pair 独立做二分类，正例在对角线。}
\figsource{Stanford CS336 2026 Lecture 17 官方可执行讲义，\texttt{images/siglip-code.png}。}
\end{figure}

\begin{knowledgebox}{读图：SigLIP 改的是归一化耦合方式}
代码同样先得到图像/文本 embeddings，但 labels 是对角线为 $+1$、非对角线为 $-1$ 的矩阵，每个 logit 独立进入 sigmoid loss。与 CLIP 的 batch-wide softmax 相比，它不要求所有设备共享一个分类分母，因此更容易分片；但负例数量、正负不平衡和 logit bias 仍会影响优化，不能把“无 global softmax”误读成“无需大规模数据或负例设计”。
\end{knowledgebox}

\begin{figure}[H]
\centering
\includegraphics[width=0.82\textwidth]{siglip-parallelism.png}
\caption{SigLIP 的 pairwise objective 更适合分片与大规模并行。}
\figsource{Stanford CS336 2026 Lecture 17 官方可执行讲义，\texttt{images/siglip-parallelism.png}。}
\end{figure}

\begin{warningbox}{对齐 embedding 不等于细粒度理解}
CLIP/SigLIP 很适合检索和 global semantics，但 spatial relations、OCR、counting 和 fine-grained grounding 需要更高分辨率 tokens、任务数据和跨层融合。把一个高 CLIP score 当成完整视觉理解会高估模型。
\end{warningbox}

\subsection{本章小结}

CLIP/SigLIP 学到可迁移视觉 encoder，解决了“图像怎样进入语言空间”的第一步。下一步是把这些视觉 features 真正注入 autoregressive LM。

